\documentclass[11pt]{article}

% --------------------------------------------------
% Fonts and Layout (LuaLaTeX)
% --------------------------------------------------
\usepackage{fontspec}
\setmainfont{Latin Modern Roman}

\usepackage[margin=1in]{geometry}
\usepackage{setspace}
\setstretch{1.15}

\usepackage{microtype}
\usepackage{amsmath}
\usepackage{hyperref}

\hypersetup{
  colorlinks=true,
  linkcolor=blue,
  urlcolor=blue
}

\title{Admissible Histories: Incongruent Neurons and the Temporal Structure of Decision-Making}
\author{Flyxion}
\date{\today}

\begin{document}
\maketitle

\begin{abstract}
Errors in decision-making are typically explained as the result of noise, insufficient evidence, or late-stage failures of control. Recent findings challenge this view. A biomimetic corticostriatal model and corresponding macaque data reveal a population of neurons—termed incongruent neurons—whose activity during the first 200 ms of a trial reliably predicts an incorrect behavioral response occurring more than one second later (Pathak et al., 2025). These neurons are stable, selective, and causally upstream of error. In this paper, we argue that such findings require a reconceptualization of error. Rather than reflecting a breakdown of computation, errors reflect the early commitment to internally coherent action trajectories that are later evaluated as incorrect. We propose that decision-making is best understood as the selection and suppression of temporally extended histories rather than as the momentary computation of correct responses. This perspective clarifies the timing of commitment, the persistence of error-related neural activity, and the success of early intervention, and it has implications for theories of agency, rationality, and cognitive control.
\end{abstract}

\section{Introduction}

Errors are among the most informative events in cognitive science. They reveal the limits of perception, the structure of decision processes, and the constraints under which cognitive systems operate. Despite their importance, errors are usually treated as secondary phenomena: deviations from an otherwise correct computation. In classical decision models, errors arise when noise overwhelms signal, when evidence accumulation terminates prematurely, or when control mechanisms fail to suppress incorrect responses (Ratcliff, 1978; Bogacz et al., 2006; Ratcliff and McKoon, 2008). Under these accounts, correct behavior reflects the intended operation of the system, while error reflects its failure.

Recent work complicates this picture. Using a biologically grounded model of corticostriatal micro-assemblies, Pathak et al. (2025) identified a class of neurons whose activity reliably predicts erroneous behavioral outcomes long before those outcomes are expressed. These “incongruent neurons” fire selectively within the first 200 ms of stimulus presentation and predict errors that occur more than one second later. Importantly, these neurons are not weak or noisy signals. They are stable across trials, selective for specific stimulus–response combinations, and causally involved in driving the incorrect action.

The temporal profile of this activity presents a challenge to standard interpretations. If an incorrect outcome can be predicted with high confidence so early in a trial, then the decisive commitment must occur well before the moment of response. The long delay between early neural activity and late behavioral execution cannot be straightforwardly interpreted as deliberation or continued decision formation. Instead, it suggests that the system has already committed to a particular course of action, and that subsequent processing reflects the unfolding of that commitment rather than its resolution.

This observation raises a fundamental question: what does it mean to make an error if the outcome is fixed before deliberation appears to conclude? One possibility is that the brain occasionally miscomputes the correct answer very quickly. Another, explored here, is that the brain computes something different altogether. Rather than computing correct actions and occasionally failing, it may generate multiple possible action trajectories and select among them under constraint. Errors, on this view, are not failures of computation but selections of trajectories that are internally coherent yet externally disfavored.

In this paper, we develop this second interpretation. We argue that incongruent neurons reveal a trajectory-based structure of decision-making in which early neural activity authorizes complete future courses of action. Errors arise when one such course—perfectly coherent given the system’s internal dynamics—is evaluated negatively by task feedback. We show how this perspective clarifies the timing of commitment, explains the persistence of error-related neurons, and makes sense of real-time intervention results. We conclude by discussing implications for theories of agency and cognitive control.

\section{Early Neural Commitment}

The most striking feature of incongruent neuron activity is its timing. In Pathak et al. (2025), these neurons fire within approximately 200 ms of stimulus onset, yet the associated behavioral response occurs more than 1,000 ms later. During this extended interval, the system continues to process information, maintain posture, and execute motor plans, but the eventual outcome is already determined.

This temporal dissociation challenges the common assumption that decision formation unfolds gradually until a threshold is crossed near the time of response. In drift-diffusion models, for example, early activity is expected to be ambiguous, reflecting partial evidence accumulation, while commitment occurs late (Ratcliff and McKoon, 2008). Incongruent neurons invert this relationship: early activity is decisive, and late activity appears epiphenomenal with respect to outcome selection.

One way to understand this pattern is to distinguish between commitment and expression. Commitment refers to the point at which the system enters a state from which only one outcome is reachable, given its internal dynamics. Expression refers to the eventual motor realization of that outcome. Incongruent neurons appear to mark the moment of commitment rather than the moment of expression.

This distinction aligns with earlier findings showing that neural activity predictive of action precedes conscious awareness of intention (Libet et al., 1983). However, incongruent neurons extend this logic to errors specifically. They show that incorrect outcomes are not the result of late slips or execution failures, but of early commitments that unfold deterministically.

\section{Error Without Failure}

A central implication of this early commitment is that error cannot be straightforwardly interpreted as computational failure. Incongruent neurons do not exhibit signatures of instability, indecision, or noise. They are selective, repeatable, and causally effective. From the system’s internal perspective, the trajectory they support is as well-formed as any trajectory leading to a correct response.

This observation motivates a distinction between internal coherence and external evaluation. A trajectory is internally coherent if it is consistent with the system’s architecture, connectivity, and learning history. It is externally correct or incorrect depending on how its outcome is evaluated by the environment. These two properties need not coincide.

Learning systems routinely preserve internally coherent behaviors that are no longer optimal. Reinforcement learning adjusts the relative strength of action pathways but rarely eliminates alternatives entirely (Sutton and Barto, 1998). As a result, disfavored behaviors remain latent and can be expressed under certain conditions. Incongruent neurons appear to be the neural signature of such latent alternatives.

Under this interpretation, errors are not deviations from admissible behavior. They are admissible behaviors that receive negative feedback. The surprise, then, is not that such behaviors exist, but that they can be identified so early and so reliably.

\section{Why Error Signals Persist}

If incongruent neurons support disfavored trajectories, why are they not eliminated by learning? The answer lies in the structure of biological learning and control. Cortical and corticostriatal circuits are sparsely connected, shaped by developmental constraints, and organized around competitive dynamics (Miller and Cohen, 2001). Learning modifies synaptic weights incrementally, amplifying rewarded pathways and attenuating unrewarded ones, but rarely deleting pathways altogether.

This incremental modification preserves flexibility. A system that eliminated all but one trajectory would be brittle, unable to adapt when task demands change. Retaining alternative trajectories allows for rapid reconfiguration at the cost of occasional error. From this perspective, incongruent neurons are not malfunctions but structural necessities.

Their persistence also explains why error rates never fall to zero, even in well-trained systems. Errors are not merely the residue of noise; they are the occasional realization of alternatives that remain dynamically available.

\section{Intervention and Control}

The strongest evidence for this trajectory-based interpretation comes from intervention. Pathak et al. (2025) showed that halting trials in real time upon detection of incongruent neuron activity significantly improves accuracy. Crucially, this intervention occurs before any incorrect behavior is expressed.

This result cannot be explained by error correction in the usual sense. There is no incorrect action to correct at the time of intervention. Instead, the intervention prevents a particular future from occurring. It acts not on a state but on a trajectory.

The success of this intervention demonstrates that error trajectories are fully specified before behavior. It also suggests a new way of thinking about cognitive control. Control may consist less in correcting actions once initiated and more in suppressing disfavored trajectories before they unfold. In everyday terms, self-control may be a matter of stopping oneself early rather than fixing mistakes later.

\section{Implications and Conclusions}

The discovery of incongruent neurons invites a reconceptualization of decision-making and error. Rather than viewing decisions as momentary computations that sometimes fail, we can view them as commitments to temporally extended courses of action. Errors arise when one such course, internally coherent but externally disfavored, is allowed to unfold.

This perspective has implications for theories of rationality, agency, and cognitive control. Rationality may consist not in computing optimal actions but in effectively suppressing suboptimal trajectories. Agency may consist not in choosing freely at the moment of action but in shaping the space of possible futures early.

More broadly, the findings suggest that cognitive systems operate over histories rather than states. Understanding behavior, therefore, requires understanding how these histories are generated, weighted, and excluded. Incongruent neurons provide a rare empirical window into this deeper structure, revealing that even our mistakes are, in a sense, chosen early.

\begin{thebibliography}{99}

\bibitem{Pathak2025}
Pathak, A., Imafuku, J., Burke, M. G., McDougle, S. R., \& Cohen, J. D. (2025).
Biomimetic model of corticostriatal micro-assemblies discovers a neural code.
\emph{Nature Communications}, 16, 7076.

\bibitem{Ratcliff1978}
Ratcliff, R. (1978).
A theory of memory retrieval.
\emph{Psychological Review}, 85(2), 59--108.

\bibitem{RatcliffMcKoon2008}
Ratcliff, R., \& McKoon, G. (2008).
The diffusion decision model.
\emph{Neural Computation}, 20(4), 873--922.

\bibitem{Bogacz2006}
Bogacz, R., Brown, E., Moehlis, J., Holmes, P., \& Cohen, J. D. (2006).
The physics of optimal decision making.
\emph{Psychological Review}, 113(4), 700--765.

\bibitem{Libet1983}
Libet, B., Gleason, C. A., Wright, E. W., \& Pearl, D. K. (1983).
Time of conscious intention to act.
\emph{Brain}, 106(3), 623--642.

\bibitem{MillerCohen2001}
Miller, E. K., \& Cohen, J. D. (2001).
An integrative theory of prefrontal cortex function.
\emph{Annual Review of Neuroscience}, 24, 167--202.

\bibitem{SuttonBarto1998}
Sutton, R. S., \& Barto, A. G. (1998).
\emph{Reinforcement Learning: An Introduction}.
MIT Press.

\end{thebibliography}

\end{document}
